%% kapa: exam=1 round=20 year=2009 period=2 form=A subject=appraisal_law subject_ko=부동산관계법규 seq=39 number=79
%% source: Q-Net 감정평가사 기출문제 (cst00309) artlSeq=5024178 file=extracted/1008/5024178/2009년 제20회 감정평가사 1차시험 기출문제/2교시.hwp
%% conversion: kordoc 4.24.51 -> tools/gen1.py (auto); review=proofread-sonnet; flags=none
\begin{question}{79}
부동산등기법령상 등기의 대상이 되는 권리에 관한 설명으로 옳지 \uline{않은} 것은?
\begin{choices}
  \item 유치권은 부동산물권이므로 등기를 할 수 있다.
  \item 건물의 일부에 대한 전세권도 등기할 수 있다.
  \item 저당권에 의하여 담보된 채권을 질권의 목적으로 하는 경우에도 등기할 수 있다.
  \item 부동산물권변동을 목적으로 하는 채권적 청구권도 가등기의 형식으로 등기할 수 있다.
  \item 부동산임차권과 환매권은 채권이지만 등기할 수 있다.
\end{choices}
\end{question}
\answer{1}
